\appendix
\section{Cline Prompts Used During Development}
\label{app:cline-prompts}

This appendix contains the archived prompts used during the development of
ServiceConnect. The prompts are presented in their recorded form.


\subsection*{MASTER CURSOR PROMPT --- Phase 0}

\begin{quote}
You are the lead full-stack engineer helping me build a production-style
college project called ServiceConnect.

ServiceConnect is a local service discovery and booking marketplace.

Customers can discover local service providers such as plumbers, mechanics,
electricians, carpenters, painters, cleaners, appliance repair technicians,
interior designers, etc.

There are three roles:

1. CUSTOMER
2. PROVIDER
3. ADMIN

Customers can search/filter providers and create service bookings.

Providers can manage their profiles, services, availability, booking requests,
and clients.

Admins manually verify providers before they become publicly visible.

IMPORTANT:
Do not build the entire application in one step.

First inspect the current workspace.

If the workspace is empty, propose and create the project structure using:

Frontend:
React + TypeScript + Vite + Tailwind CSS

Backend:
Node.js + Express + TypeScript

Database:
PostgreSQL

Authentication:
Secure role-based authentication

Before making large changes:
1. Explain the proposed architecture.
2. Create a clean folder structure.
3. Create README.md.
4. Create .env.example.
5. Create docs/BUILD_LOG.md.
6. Create docs/AI_DECISION_LOG.md.
7. Create docs/SECURITY.md.
8. Create docs/DEPLOYMENT.md.

Do not put secrets into source code.

Use modular architecture.

Do not use fake security controls implemented only in the frontend.

Backend authorization must enforce CUSTOMER, PROVIDER, and ADMIN permissions.

Start with project scaffolding only.

After scaffolding, stop and tell me:
- what files were created
- how to install dependencies
- how to start the frontend
- how to start the backend
- what remains to implement

Do not continue automatically into advanced features.
\end{quote}

\subsection*{PHASE 1 --- Project Scaffolding Prompt}

\begin{quote}
Now implement the initial ServiceConnect project scaffolding.

Requirements:

Frontend:
- React
- TypeScript
- Vite
- Tailwind CSS
- React Router
- Lucide icons

Backend:
- Node.js
- Express
- TypeScript
- centralized error handling
- environment configuration
- API routing structure

Create:
client/
server/
docs/
database/

Add a health endpoint:

GET /api/health

Response:

{
  "status": "ok",
  "service": "ServiceConnect API"
}

Create a clean frontend landing page placeholder.

Create a responsive navigation component.

Do not implement authentication yet.

Run/build/test the project if possible.

Fix all errors before finishing.

At the end, summarize:
1. files created
2. commands used
3. tests performed
4. issues encountered
5. decisions made
\end{quote}

\subsection*{PHASE 2 --- Database Prompt}

\begin{quote}
Now implement the ServiceConnect database layer.

Use PostgreSQL.

Create the following entities:

User
ProviderProfile
ServiceCategory
Service
Booking
Review
ProviderImage
AdminActionLog

Implement proper relationships and constraints.

User roles:
CUSTOMER
PROVIDER
ADMIN

Provider verification statuses:
PENDING
APPROVED
REJECTED
SUSPENDED

Booking statuses:
PENDING
ACCEPTED
REJECTED
CANCELLED
IN_PROGRESS
COMPLETED

Add:
- primary keys
- foreign keys
- indexes where appropriate
- timestamps
- useful uniqueness constraints

Important:
Passwords must never be stored in plaintext.

Create migrations.

Create seed data for development.

Use fictional/demo information only.

Add seed categories and at least 8 fictional providers.

Only APPROVED providers should be considered public.

Document database decisions in docs/AI_DECISION_LOG.md.

Do not continue to frontend features until the database layer works.
\end{quote}

\subsection*{PHASE 3 --- Authentication Prompt}

\begin{quote}
Implement secure authentication for ServiceConnect.

Requirements:

- customer registration
- provider registration
- login
- logout
- current-user endpoint
- password hashing
- authentication middleware
- role-based authorization middleware

Roles:
CUSTOMER
PROVIDER
ADMIN

Security requirements:

- never store plaintext passwords
- validate input
- do not expose password hashes
- do not trust role information supplied by the frontend
- protect private routes
- protect admin routes
- use environment variables for secrets
- implement appropriate token/session expiration
- return safe error messages
- do not reveal whether sensitive credentials exist unnecessarily

Create:

POST /api/auth/register
POST /api/auth/login
POST /api/auth/logout
GET /api/auth/me

Create reusable middleware such as:

requireAuth
requireRole

Add tests for:
- successful registration
- duplicate email
- invalid login
- protected endpoint without authentication
- customer accessing admin endpoint
- provider accessing admin endpoint

Document security decisions.
\end{quote}

\subsection*{PHASE 4 --- Provider Registration + Verification Prompt}

\begin{quote}
Implement the provider onboarding and verification workflow.

Flow:

Provider registers
→ profile created
→ status PENDING
→ provider is NOT publicly visible
→ admin reviews provider
→ admin approves/rejects
→ approved provider becomes publicly searchable

Provider profile should support:

- business name
- description
- phone
- location
- service area
- experience
- profile image
- cover image

Implement:

GET /api/providers/me
PATCH /api/providers/me

Admin:

GET /api/admin/providers/pending
GET /api/admin/providers/:id
PATCH /api/admin/providers/:id/approve
PATCH /api/admin/providers/:id/reject
PATCH /api/admin/providers/:id/suspend

Reject should support a reason.

Every admin decision must create an AdminActionLog.

Security:
- provider cannot approve themselves
- provider cannot modify verificationStatus
- customer cannot access provider admin endpoints
- provider A cannot modify provider B
- only admin can change verification status

Add tests for these authorization rules.
\end{quote}

\subsection*{PHASE 5 --- Service Management Prompt}

\begin{quote}
Implement provider service management.

Providers can:

- create services
- edit services
- deactivate services
- view their services

A service contains:

- category
- name
- description
- priceFrom
- priceTo
- estimated duration
- active status

Endpoints:

GET /api/providers/me/services
POST /api/providers/me/services
PATCH /api/providers/me/services/:id
DELETE /api/providers/me/services/:id

Important:
A provider may only modify their own services.

A provider must never be able to modify another provider's service
by changing an ID in the request URL.

Implement ownership checks server-side.

Add validation for:
- price
- service name
- description
- duration

Add tests for IDOR/ownership vulnerabilities.
\end{quote}

\subsection*{PHASE 6 --- Customer Search Prompt}

\begin{quote}
Now implement the public provider search system.

Customers should be able to search approved providers.

Filters:

- keyword
- category
- location
- minPrice
- maxPrice
- rating
- availability

Only providers with verificationStatus=APPROVED may appear.

Create:

GET /api/providers

Support query parameters.

Validate all query parameters.

Prevent SQL injection by using parameterized queries or ORM-safe queries.

Implement pagination.

Implement sorting:

- rating
- price
- newest

Return only information appropriate for public users.

Do not expose:
- password hashes
- private admin information
- internal verification notes
- sensitive personal information

Add API tests for:
- filtering
- pagination
- unauthorized data exposure
- pending provider visibility
- rejected provider visibility
\end{quote}

\subsection*{PHASE 7 --- Frontend Home Page Prompt}

\begin{quote}
Now build the ServiceConnect customer-facing home page.

Design goals:

- modern
- aesthetic
- professional
- responsive
- visually engaging
- easy to understand

Hero text:

Find reliable local services.
Book the right professional for the job.

Hero search:

What service do you need?
Where?

CTA:
Search Services

Add sections:

1. Popular Services
2. How ServiceConnect Works
3. Featured Providers
4. Why ServiceConnect
5. Call-to-action for providers
6. Footer

Popular categories:

Plumbing
Electrical
Carpentry
Painting
Mechanics
Cleaning
AC Repair
Interior Design

Use cards and icons.

Add subtle animations only where useful.

Make the page responsive.

Do not use generic placeholder dashboard styling.

Use a consistent ServiceConnect design system.
\end{quote}

\subsection*{PHASE 8 --- Search UI Prompt}

\begin{quote}
Build the provider search page.

Requirements:

- search bar
- location input
- category selection
- price range
- rating filter
- availability filter
- sort dropdown
- result count
- provider cards
- loading state
- empty state
- error state

Provider card should show:

- image
- business name
- verified badge
- category
- location
- rating
- review count
- starting price
- availability
- View Profile
- Book Service

On mobile:
- filters should become a drawer/modal
- cards should stack
- search controls should remain usable

Connect the UI to the actual backend API.

Do not use hard-coded provider results once API integration is available.
\end{quote}

\subsection*{PHASE 9 --- Provider Profile Prompt}

\begin{quote}
Build the provider profile page.

Show:

- cover image
- profile image
- business name
- verified badge
- rating
- review count
- location
- description
- services
- price range
- experience
- availability
- image gallery
- reviews
- booking CTA

Only approved providers should be publicly accessible.

If a provider is pending/rejected/suspended, do not expose the normal
public profile to customers.

Use a professional marketplace layout.

Make it responsive.

Add loading/error states.
\end{quote}

\subsection*{PHASE 10 --- Booking Backend Prompt}

\begin{quote}
Implement the ServiceConnect booking system.

Customer booking flow:

1. select provider
2. select service
3. describe problem
4. select date
5. select time
6. provide service address
7. add notes
8. submit request

Booking statuses:

PENDING
ACCEPTED
REJECTED
CANCELLED
IN_PROGRESS
COMPLETED

Endpoints:

POST /api/bookings
GET /api/bookings/my
GET /api/bookings/:id
PATCH /api/bookings/:id/cancel

Provider:

GET /api/provider/bookings
PATCH /api/provider/bookings/:id/status

Security:

Customer can only create/view/cancel their own bookings.

Provider can only view/update bookings assigned to them.

Provider cannot modify another provider's booking.

Customer cannot change booking status to ACCEPTED or COMPLETED.

Validate dates and required fields.

Prevent invalid state transitions.

Add tests for ownership and authorization.
\end{quote}

\subsection*{PHASE 11 --- Booking UI Prompt}

\begin{quote}
Build the complete customer booking experience.

Create a clean multi-step booking interface.

Step 1:
Select service

Step 2:
Describe problem

Step 3:
Date and time

Step 4:
Address

Step 5:
Review

Step 6:
Confirmation

After booking:

Show:

Booking submitted successfully.

Provider:
ABC Plumbing

Service:
Leak Repair

Date:
30 September 2026

Status:
Pending

Add a View Booking button.

Use validation and clear error messages.

Do not allow submission when required fields are invalid.
\end{quote}

\subsection*{PHASE 12 --- Customer Dashboard Prompt}

\begin{quote}
Build the customer dashboard.

Sections:

- overview
- upcoming booking
- active requests
- booking history
- saved providers
- profile

Show useful cards:

Upcoming Booking
Pending Requests
Completed Services

Create a responsive sidebar/navigation.

Add loading, empty, and error states.

Ensure customers only receive their own data.
\end{quote}

\subsection*{PHASE 13 --- Provider Dashboard Prompt}

\begin{quote}
Build the provider dashboard.

Show:

- today's bookings
- pending requests
- completed jobs
- total clients
- upcoming schedule

Sections:

- Overview
- Booking Requests
- Upcoming Bookings
- Clients
- Services
- Availability
- Business Profile
- Reviews
- Settings

Provider booking actions:

Accept
Reject
Start Job
Complete Job

Validate status transitions.

Do not allow providers to manipulate bookings belonging to other providers.

Make the dashboard visually polished and responsive.
\end{quote}

\subsection*{PHASE 14 --- Admin Dashboard Prompt}

\begin{quote}
Build the ServiceConnect admin dashboard.

Admin dashboard should include:

- total customers
- total providers
- pending providers
- total bookings
- active bookings

Create provider verification screen.

Each pending provider should display:

- business name
- provider name
- category
- location
- contact
- description
- submitted images
- services
- submitted date
- status

Actions:

Approve
Reject
Suspend
View Details

Reject action should request a reason.

Create activity log page.

Admin UI should be clearly separated from customer/provider interfaces.

SECURITY:
Do not rely on frontend route hiding.

Every admin API must verify:
1. authentication
2. ADMIN role
\end{quote}

\subsection*{PHASE 15 --- Image Upload Prompt}

\begin{quote}
Implement secure provider image uploads.

Providers can upload:

- profile image
- cover image
- business/gallery images

Security requirements:

- validate MIME type
- validate extension
- limit file size
- generate safe filenames
- do not trust original filename
- prevent executable uploads
- store files outside sensitive source directories
- return safe URLs
- handle upload failures
- restrict who can upload
- restrict who can delete images

Document upload security in docs/SECURITY.md.

Add tests for invalid file types and oversized uploads.
\end{quote}

\subsection*{PHASE 16 --- Reviews Prompt}

\begin{quote}
Implement reviews.

Only customers with a COMPLETED booking should be allowed to review
the corresponding provider/service.

Review fields:

rating 1-5
comment

Prevent:

- duplicate reviews for the same booking
- reviews without completed booking
- users reviewing unrelated providers
- providers reviewing themselves
- unauthenticated reviews

Show provider rating and review count.

Validate review input.
\end{quote}

\subsection*{PHASE 17 --- UI Polish Prompt}

\begin{quote}
Now perform a complete UI/UX refinement pass.

Do not add major backend features.

Inspect the entire application.

Improve:

- spacing
- typography
- buttons
- cards
- navigation
- forms
- responsive layouts
- loading states
- empty states
- error states
- toast messages
- status badges
- icons
- accessibility
- visual consistency

Make ServiceConnect look like a real commercial local-services platform.

Avoid:
- excessive gradients
- excessive animations
- clutter
- inconsistent colors
- inconsistent border radius
- generic default styles

Check every page on desktop and mobile.

Fix obvious visual bugs.
\end{quote}

\subsection*{PHASE 18 --- Security Audit Prompt}

\begin{quote}
Act as a security engineer.

Do a security review of the entire ServiceConnect codebase.

Do not simply tell me that it is secure.

Inspect the actual implementation.

Check for:

1. Authentication weaknesses
2. Authorization weaknesses
3. IDOR
4. Broken role checks
5. SQL injection
6. XSS
7. CSRF where relevant
8. Unsafe file uploads
9. Weak password handling
10. Sensitive data exposure
11. Hard-coded secrets
12. Excessive CORS
13. Missing rate limiting
14. Insecure error messages
15. Missing input validation
16. Booking manipulation
17. Provider/customer data leakage
18. Admin endpoint exposure
19. Insecure direct object access
20. Dependency vulnerabilities
21. Insecure environment configuration
22. Missing ownership checks

For each finding provide:

- Finding
- Location/file
- Severity
- Explanation
- Exploit scenario
- Recommended fix

Do not modify code yet.

Save the review to:

docs/SECURITY_REVIEW.md
\end{quote}

\subsection*{PHASE 19 --- Security Fix Prompt}

\begin{quote}
Using docs/SECURITY_REVIEW.md, fix the identified security vulnerabilities.

Prioritize:

CRITICAL
HIGH
MEDIUM
LOW

For every fix:

1. Explain the vulnerability.
2. Explain the fix.
3. Modify the minimum required code.
4. Add a regression test where possible.
5. Run tests.
6. Verify the vulnerability is no longer reproducible.

Do not introduce unrelated changes.

Update docs/SECURITY.md.

Record the changes in docs/AI_DECISION_LOG.md.
\end{quote}

\subsection*{PHASE 20 --- Testing Prompt}

\begin{quote}
Create a complete testing pass.

Test:

Authentication
Authorization
Customer flows
Provider flows
Admin flows
Provider verification
Service CRUD
Search
Filters
Bookings
Booking state transitions
Reviews
Image uploads
API validation

Include negative tests.

Important negative tests:

- unauthenticated request
- customer accessing provider endpoint
- provider accessing admin endpoint
- provider modifying another provider
- customer accessing another customer's booking
- provider accessing another provider's booking
- changing verification status as non-admin
- invalid booking status transition
- invalid file upload
- malformed input
- invalid price
- invalid date

Fix failing tests.

Provide final test command documentation in README.md.
\end{quote}

\subsection*{PHASE 21 --- Deployment Prompt}

\begin{quote}
Prepare ServiceConnect for deployment.

Create/update:

README.md
docs/DEPLOYMENT.md
.env.example

Document:

Frontend build command
Backend build command
Backend start command
Database migration
Database seed
Required environment variables
CORS configuration
Production API URL
Frontend API URL
File storage requirements

Verify:

- no secrets in Git
- .env is ignored
- production configuration is separated
- debug logging is disabled where appropriate
- CORS is configured
- database connection uses environment variables

If Docker is used, create/update:

Dockerfile
docker-compose.yml

Do not deploy yet unless explicitly requested.
\end{quote}

\subsection*{PHASE 22 --- Final Code Review Prompt}

\begin{quote}
Perform a final engineering review of ServiceConnect.

Review:

Architecture
Code organization
Frontend quality
Backend quality
Database design
Security
Authentication
Authorization
Validation
API consistency
Error handling
Testing
Responsive design
Accessibility
Deployment readiness
Documentation

Identify:

- technical debt
- duplicated code
- unused files
- dead code
- insecure patterns
- missing tests
- poor naming
- performance issues

Fix only issues that are appropriate for this project.

Do not rewrite the entire application unnecessarily.

At the end produce:

FINAL_REVIEW.md

with:

1. Completed features
2. Remaining limitations
3. Security controls
4. Tests
5. Deployment readiness
6. Known limitations
7. Future improvements
\end{quote}

\subsection*{Phase 23 --- AI Build Log Finalization Prompt}

\begin{quote}
Now help me prepare the AIEngineering assignment documentation.

Review the Git history, docs, and project changes.

Create/update:

docs/BUILD_LOG.md

For each major development phase record:

- Date
- Objective
- Prompt used
- AI-generated solution
- What was accepted
- What was rejected
- What was modified manually
- Why the decision was made
- Files changed
- Tests performed
- Result

Create:

docs/AI_DECISION_LOG.md

Include important architectural and UI decisions.

Do not invent decisions that were not actually made.

If Git history does not provide enough information, clearly mark
the item as "Not recorded" rather than fabricating details.
\end{quote}

\subsection*{Phase 24 --- Peer Audit Readiness Prompt}

\begin{quote}
Prepare ServiceConnect for another student group to audit.

Create:

docs/PEER_AUDIT_GUIDE.md

Include:

1. Application overview
2. Technology stack
3. How to run locally
4. Test accounts
5. Role descriptions
6. Important endpoints
7. Security controls
8. Known limitations
9. Areas suitable for security testing
10. Deployment configuration

Do not reveal secrets.

Use only development/demo credentials.

Also create:

docs/AUDIT_RESPONSE_TEMPLATE.md

with sections:

Finding
Severity
Validity
Affected Component
Root Cause
Fix
Testing
Status
\end{quote}

\subsection*{Prompt for Fixing Any Cursor Error}

\begin{quote}
Do not rewrite the whole project.

First inspect the error and trace it to the root cause.

Explain:
1. What caused the error.
2. Which file is responsible.
3. Why the current implementation fails.
4. The smallest safe fix.

Then implement the fix.

After fixing:
- run the relevant test/build
- verify the original functionality
- check that the fix did not break unrelated features

Do not change unrelated code.
\end{quote}

\subsection*{Prompt for UI Problems}

\begin{quote}
Inspect the current UI implementation and fix the following issue:

[DESCRIBE ISSUE]

Do not redesign the entire website.

Keep the existing ServiceConnect design system.

Fix:
- responsiveness
- spacing
- alignment
- typography
- component sizing
- overflow
- accessibility

Verify desktop and mobile layouts.

Explain what changed and why.
\end{quote}

\subsection*{Prompt for AI Hallucinated/Bad Code}

\begin{quote}
Review your previous implementation critically.

Do not assume the previous code is correct.

Check the actual repository files.

If your previous implementation conflicts with the project requirements,
identify the conflict.

Do not defend the previous implementation.

Correct it using the project specification.

Before changing code:
- identify the problem
- explain the correct behavior
- implement the smallest appropriate fix
- test it
\end{quote}

\subsection*{Correction 1 --- Indian/Goa localization and INR price conversion}

\begin{quote}
Update ServiceConnect for the intended Indian/Goa context.
Requirements:
- use INR / ₹ instead of USD
- use Indian number/date formatting
- replace US-style demo locations with Goa/India data
- preserve the existing architecture
- preserve authentication and authorization
- preserve provider verification
- preserve customer/provider/admin roles
- ensure provider registration works
- ensure role-based routing works
- ensure only approved providers are public
- update relevant tests
- run server/client tests
- run typechecks
- run production build
Do not make unrelated changes.
\end{quote}

\subsection*{Correction 2 --- Provider dashboard crash}

\begin{quote}
Inspect the provider dashboard crash and trace it to the actual API
response.
Do not rewrite the dashboard.
Correct the frontend/API contract mismatch.
Use the actual response shape.
Update the relevant type/API code and add or update regression tests.
Run:
- server tests
- client tests
- typechecks
- production build
Do not change unrelated functionality.
\end{quote}

\subsection*{Correction 3 --- Provider business page service-shape crash}

\begin{quote}
Inspect the actual owner-service API response.
Correct the ProviderService type and provider business page to use the
real flat response.
Do not alter the public service API unnecessarily.
Add regression tests for the provider business page.
Run:
- server tests
- client tests
- typechecks
- production build
Do not change unrelated code.
\end{quote}

\subsection*{Correction 4 --- Custom service categories}

\begin{quote}
Allow providers to create a custom service category while preserving the
existing predefined categories.
A service may specify exactly one of:
- categoryId
- categoryName
Normalize category names by:
- trimming
- collapsing repeated whitespace
- normalizing the slug
Examples:
CCTV Installation
cctv installation
CCTV Installation
must resolve to the same category:
cctv-installation
If a matching category already exists, reuse it.
If it does not exist, create it.
Handle concurrent category creation safely:
- attempt insertion
- handle unique constraint race
- re-read the existing category
Custom categories should receive a high sort order so predefined
categories remain first.
Frontend:
- use a searchable/free-text combobox
- selecting an existing category pins its ID
- typing a custom category clears the selected ID
- show filtered suggestions
Add tests for:
- normalization
- duplicate categories
- concurrent creation
- existing category reuse
- custom category creation
Run the full test/typecheck/build suite.
\end{quote}

\subsection*{Correction 5 --- Secure provider business image uploads}

\begin{quote}
Implement secure provider business image uploads.
Providers should be able to:
- upload business/gallery images
- view their uploaded images
- delete their own images
Endpoints:
GET /api/providers/me/images
POST /api/providers/me/images
DELETE /api/providers/me/images/:id
Security:
- provider authentication required
- provider ownership enforced
- MIME type validation
- extension validation
- file size limit
- safe UUID filenames
- no executable uploads
- upload failures handled safely
- safe public URLs
Use:
server/uploads/providers
Serve through:
/uploads/providers
Ignore uploaded files in Git.
Add backend and frontend tests.
Run:
- server tests
- client tests
- typechecks
- production build
Do not change unrelated features.
\end{quote}

\subsection*{Correction 6 --- Business images, admin review, service images, dummy-service cleanup}

\begin{quote}
Expand the existing image functionality.
1. Customer business images
Customers must be able to see approved provider business/gallery images
on the public provider profile.

2. Admin review images
When an admin reviews a provider, the admin must see the provider's
business images.
Admin image access is read-only.
Do not add admin image mutation endpoints.

3. Service-specific images
Providers must be able to upload images for individual services.
Add a service_images table with:
- id
- service_id
- url
- alt_text
- sort_order
- created_at
Use:
- foreign key to services
- cascade delete
- URL validation
- unique service/image URL
- useful indexing
Endpoints:
GET /api/providers/me/services/:id/images
POST /api/providers/me/services/:id/images
DELETE /api/providers/me/services/:id/images/:imageId
Only the owning provider may manage the images.

4. Provider UI
Add:
- business image management
- service image management
For a newly created service without an ID, do not attempt to upload
service images until the service has been created.

5. Customer UI
Display service-specific images on the service listing/profile.
If a service has no images, omit the image section.

6. Dummy/default services
Remove the unwanted seeded/demo services that were not approved.
Preserve real user-created services.
Remove associated dummy data only where required by the database
relationships.
Do not recreate the removed dummy services when seed is run again.

7. Testing
Add tests for:
- business image access
- business image deletion
- service image upload
- service image deletion
- ownership
- admin read-only review
- customer visibility
- invalid files
- missing images
IMPORTANT:
Do not modify:
server/src/tests/search.test.ts
Do not change the application search implementation merely to make the
search test pass.
Run:
- complete server test suite
- complete client test suite
- server typecheck
- client typecheck
- root typecheck
- production build
After implementation, manually verify these flows in the browser:
A. Provider uploads a business image.
 Customer can see it on the public provider page.
B. Provider uploads a service image.
 Customer can see it on the corresponding service.
C. Admin opens provider review details.
 Business images are visible.
D. Provider deletes a service image.
 Customer no longer sees it.
E. Provider deletes a business image.
 Customer no longer sees it.
\end{quote}

\subsection*{Correction 7 --- Final image-work test-fixture correction}

\begin{quote}
Manual browser testing of the image features passed.

Do not modify server/src/tests/search.test.ts.

Fix ONLY the newly exposed keyword-search test failure by correcting its test fixture/setup so that it selects an APPROVED/public provider or otherwise uses an appropriate public-search fixture.

Do NOT change the application search implementation to make the test pass.

Leave the pre-existing price-sort failure unchanged.

Then run the complete server test suite, client test suite, both typechecks, and client production build.

Do not commit or push.
\end{quote}

% End of Cline prompt appendix.
